% ch9_d (书页 413-422 / p428-p437)
%--C-TAIL--

考虑一个位于 Schwarzschild 黑洞之外、半径 $r_{1} > 2GM$ 处的静态观者。这样的观者沿着类时 Killing 矢量 $K = \partial_{t}$ 的轨道运动。在第 6 章中我们已经证明，对 Schwarzschild 时空中的静态观者，红移因子 $V = \sqrt{-K_{\mu}K^{\mu}}$ 由下式给出

\begin{equation*}
V = \sqrt{1 - \frac{2GM}{r}},
\tag{9.169}
\end{equation*}

相应的加速度大小则由下式给出：

\begin{equation*}
a = \frac{GM}{r\sqrt{r - 2GM}}.
\tag{9.170}
\end{equation*}

对于非常靠近事件视界的观者，$r_{1} - 2GM \ll 2GM$，这一加速度与 Schwarzschild 半径所标定的尺度相比会变得非常大，

\begin{equation*}
a_{1} \gg \frac{1}{2GM}.
\tag{9.171}
\end{equation*}

Schwarzschild 半径又标定了视界附近时空的曲率半径。因此，在由 $a_{1}^{-1} \ll 2GM$ 所设定的长度和时间尺度上观察，时空看上去基本上是平直的。让我们作一个至关重要的假设：在黑洞附近\emph{自由下落}的观者看来，某个标量场 $\phi$ 的量子态就像 Minkowski 真空（不含任何粒子）。这个假设是合理的，因为事件视界并不是一种局域的屏障；自由下落的观者在穿过视界时看不到任何异常发生。于是，这位静态观者看起来就跟平直时空中的常加速度观者一模一样，将探测到温度为 $T_{1} = a_{1}/2\pi$ 的安鲁辐射。

现在考虑位于无限远处——或者至少位于与 $2GM$ 相比大得多的距离 $r_{2}$ 处——的静态观者。在这种情形下，无论如何都谈不上能在 $a_{2}^{-1} \gg 2GM$ 的时间尺度上忽略时空曲率，因此没有理由预期他们会看到温度为 $a_{2}/2\pi$ 的辐射，这里 $a_{2}$ 在 $r_{2}$ 处取值。但是在视界附近观测到的辐射会带着恰当的红移传播到无限远。我们可以套用上一节末尾用过的论证，来确定这样的观者应当会看到什么；他们应当探测到红移至如下温度的热辐射：

\begin{equation*}
T_{2} = \frac{V_{1}}{V_{2}} T_{1} = \frac{V_{1}}{V_{2}} \frac{a}{2\pi}.
\tag{9.172}
\end{equation*}

在无限远处 $V_{2} \to 1$，于是观测到的温度为

\begin{equation*}
T = \lim_{r_{1} \to 2GM} \frac{V_{1} a_{1}}{2\pi} = \frac{\kappa}{2\pi},
\tag{9.173}
\end{equation*}

其中 $\kappa = \lim(Va)$ 是表面引力；对 Schwarzschild 黑洞，$\kappa = 1/4GM$。与平直时空中的加速观者不同，Schwarzschild 时空中的静态 Killing 矢量在无限远处具有有限的模长，因此视界附近的辐射只红移到一个有限值，而不会一直红移到零。远离黑洞的观者因此会看到黑洞发出的热辐射流，其温度正比于黑洞的表面引力。这就是著名的\textbf{霍金效应}（Hawking effect），这种辐射本身则被称为霍金辐射。

这种对霍金效应的推导虽然取巧，却没有任何不诚实之处。特别是，与加速度的联系清楚地说明了为什么温度正比于黑洞的表面引力（这一结论对更一般的黑洞同样成立，而不只是 Schwarzschild 黑洞）。不过，我们需要对我们所作的那个假设有清醒的认识：视界附近的真空态在自由下落的观者看来是非奇异的。用术语来说，就是取重整化能动张量在视界处为有限，或者等价地，两点函数满足 Hadamard 条件 (9.123)。

通过考察最大延拓 Schwarzschild 几何中各种可能的真空态，这个假设的含义会变得更加清楚。这些态对于由引力坍缩形成的现实黑洞来说不一定具有物理上的相关性，但在这种理想化情形中出现的各种可能性，能给现实世界带来有益的教益。我们将只描述这些态，而不定量地刻画它们，也不推导它们的任何性质；更多细节可参见上文所引的文献。

在寻找真空态时，我们也许会先去寻找一个在整个时空中都正则的态［在 Hadamard 意义下，式 (9.123)］。对最大延拓的 Schwarzschild 时空，Hartle 和 Hawking 找到了这样一个态，我们称之为\textbf{Hartle--Hawking 真空}（Hartle--Hawking vacuum）；的确，这是唯一的既处处正则、又在代表无限远处时间平移的 Schwarzschild Killing 矢量 $\partial_{t}$ 作用下不变的真空态。特别是，回忆图 5.16 所示的 Schwarzschild 共形图，Hartle--Hawking 真空在 $r = 2GM$ 处的过去和未来事件视界 $H^{\pm}$ 上正则，在过去和未来类光无限远 $\mathscr{I}^{\pm}$ 上也正则。根据上文对静态观者的讨论，我们应当会预期 Hartle--Hawking 真空中有热辐射从黑洞发出，而事实证明确实如此。然而，对这个态的仔细考察揭示出，还有一份等量的热辐射流从过去类光无限远（$\mathscr{I}^{-}$）流向黑洞；换句话说，它描述的是一个与其环境处于热平衡的黑洞。我们并不会用它来模拟我们宇宙中现实的黑洞。另一种真空更类似于经由引力坍缩形成的黑洞的真空态，这就是\textbf{Unruh 真空}（Unruh vacuum），它在 $H^{+}$ 上非奇异（因此预言了出射的霍金辐射），但没有来自 $\mathscr{I}^{-}$ 的入射辐射。Unruh 真空在 Schwarzschild 的过去视界 $H^{-}$ 上是奇异的；如果我们只是把它当作现实黑洞的模型，这并不会困扰我们，因为如图 5.17 那样经历引力坍缩的时空不会有白洞，也没有任何过去视界。最后，我们也许会去寻找这样一种真空态：既没有粒子进入黑洞，也没有粒子逃往无限远；换句话说，$\mathscr{I}^{\pm}$ 上通量为零。确实存在这样一个态，称为\textbf{Boulware 真空}（Boulware vacuum）。这样一种态的存在，似乎与我们从安鲁效应出发对霍金效应所作的论证相冲突，但仔细的分析表明，Boulware 真空在 $H^{-}$ 和 $H^{+}$ 上都是奇异的。于是，认为真空在视界附近的自由下落观者看来是正则的这一假设，在这个态中被违反了。

因此，对永恒 Schwarzschild 度规中真空态的仔细考察，与我们基于安鲁效应的推理是一致的；在 $H^{+}$ 上正则的态预言了预期形式的霍金辐射。注意，事件视界的存在对这一论证至关重要；如果没有这样的视界，态必须在视界上正则的要求就没有任何约束力。以中子星为例：其半径可能接近 Schwarzschild 半径，但其时空没有任何视界。中子星不发出任何霍金辐射。理解这一点的一个途径是认识到，静态中子星度规具有一个处处类时的 Killing 矢量，可以用它来定义贯穿整个时空延伸、并且在无限远处与 Minkowski 模式相匹配的正频模式。这样得到的真空态实际上会类似于 Boulware 真空，在 $\mathscr{I}^{\pm}$ 上没有通量；而完整的 Boulware 真空在视界上奇异这一点，在中子星情形中不会困扰我们，因为那里根本不存在任何视界。

为了绝对确信我们为现实的黑洞正确选择了真空态，我们应当考虑在一个过去近乎 Minkowski 时空、未来为 Schwarzschild 时空的时空中发生的引力坍缩，如图 5.17 所示。如果真空在 $\mathscr{I}^{-}$ 上取标准的 Minkowski 形式，我们就可以追问这些模式如何穿过坍缩几何传播到 $\mathscr{I}^{+}$，这就定义了一个如式 (9.45) 中的 S 矩阵，用以确定渐近观者会看到什么。Hawking 最初发现黑洞辐射时所做的正是这件事；计算涉及一些烦琐的代数，但基本上是直截了当的，得到的温度与我们上面推导的相同。

当然，从完整的计算中我们能学到的不仅仅是黑体温度；我们会想问，比如说，当发射辐射的波长与 Schwarzschild 半径相当时会发生什么——在这种情形下我们的近似显然失效了。如果我们仔细研究任意种类的粒子从任何一种黑洞（也就是同时容许电荷和自旋）的发射，就会发现发射辐射的谱具有如下形式

\begin{equation*}
\langle \hat{n}_{\omega} \rangle = \frac{\Gamma(\omega)}{e^{2\pi(\omega - \mu)/\kappa} \pm 1}.
\tag{9.174}
\end{equation*}

这里的 $\kappa$ 当然就是表面引力。参数 $\mu$ 是一个化学势（chemical potential），刻画黑洞丢掉其守恒量子数的倾向；带电黑洞优先发射与黑洞电荷同号的粒子，而转动黑洞优先发射与黑洞角动量同号的粒子。因此，霍金辐射趋向于把黑洞带回 Schwarzschild 状态。$\Gamma(\omega)$ 是一个灰体因子（greybody factor），可以把它想象成波包被引力场反散射回黑洞而产生的。在高频极限下波长非常小，反散射可以忽略；在极低的频率下，波长变得大于 Schwarzschild 半径，反散射就变得重要起来。虽然灰体因子的解析表达式很难导出，但在频率很大和很小的极限情形下，标量场的灰体因子满足

\begin{align*}
\Gamma(\omega) &\to 1, \qquad \omega \gg \frac{1}{GM} \\
\Gamma(\omega) &\to \frac{A}{4\pi} \omega^{2}, \qquad \omega \ll \frac{1}{GM},
\tag{9.175}
\end{align*}

其中 $A$ 是黑洞的面积。

从经典广义相对论的观点来看，黑洞发出热辐射的发现当然是令人惊讶的，因为我们曾经强调过，从事件视界内部的点逃逸到无限远是不可能的。理解眼下情形的一种形象方式，是用 Feynman 图来思考真空涨落：涨落由不断地出现又消失的虚粒子/反粒子对来表示。这幅图像也很有帮助，例如有助于理解诸如 Lamb 位移（Lamb shift）这样的观测现象——原子光谱会受到光子与虚电子/正电子对相互作用的影响。通常情况下，这些粒子对总会湮灭，它们的效应只是间接的，通过与虚粒子相耦合的过程的重整化体现出来。然而，在存在事件视界时，虚粒子对中的一个成员偶尔会落入黑洞，而它的同伴则逃逸到无限远，如图 9.2 所描绘。在这幅图像中，我们观测为霍金辐射的正是这些逃逸的虚粒子。虚粒子对的总能量必须相加为零，但从无限远来看，落入的粒子可以带有负能量，因为渐近类时的 Killing 矢量在视界内部是类空的。这幅图像有些不正规，但它为正在发生的事情提供了一种有用的启发。

%FIG{9.2}: {真空涨落偶尔会导致粒子/反粒子对中的一个落入事件视界，而另一个则作为霍金辐射逃逸到无限远。}

一旦知道了黑洞温度的公式，我们就能够确定黑洞参数与热力学变量之间各关系式中的比例常数，如式 (6.118) 所列。霍金辐射可以说最终圆满成就了黑洞力学与热力学的联姻；稳态黑洞的行为就如同能量 $E = M$、温度 $T = \kappa/2\pi$、熵 $S = A/4G$ 的热平衡物体。这确实是极其巨大的熵。对宇宙中的物质场而言，熵近似等于相对论性粒子的数目；在一个 Hubble 半径之内，这个数目算出来是

\begin{equation*}
S_{\mathrm{M}} \sim 10^{88}.
\tag{9.176}
\end{equation*}

另一方面，黑洞的熵是以 Planck 单位量度的视界面积（记住，我们从头到尾都在取 $\hbar = 1$）。我们可以换算成天体物理单位，得到

\begin{equation*}
S_{\mathrm{BH}} \sim 10^{90} \left( \frac{M}{10^{6} M_{\odot}} \right)^{2}.
\tag{9.177}
\end{equation*}

这样，单个百万太阳质量的黑洞（例如在我们自己银河系的中心以及许多其他星系的中心就能找到）所具有的熵，就比可见宇宙中全部物质的熵还要多。宇宙的总熵远比我们本来能把它做到的要小——只要把更多的质量投入黑洞就行。（当宇宙学家说熵 $S_{\mathrm{M}}$ 很大时，他们的意思是：在一个曲率半径之内竟然存在这么多熵，这很令人惊讶。）我们之所以处在这样一个低熵状态，其原因想必与初始条件有关，或许还与暴胀有关。

回到黑洞力学，我们看到一个难题：宏观黑洞的熵会非常巨大，但从统计力学的观点来看，熵本应量度可及态数目的对数。经典黑洞由少数几个参数（质量、电荷和自旋）确定，因此很难知道那些态可能是什么。尽管如此，我们可以采取这样的态度：这个矛盾其实无关紧要，因为关于黑洞状态的任何信息想必都藏在事件视界背后。

把量子力学包括进来，使这个难题变得更糟而非更好，因为黑洞不仅会辐射，而且会蒸发。当我们开始研究弯曲时空中的 QFT 时，我们定下的规则之一就是：假定背景度规固定不变，不去操心量子场自身的能动张量会产生什么影响。尽管如此，即使在量子力学中我们也有能量守恒（例如，表现为渐近平直时空中守恒的 ADM 质量）。因此，当霍金辐射逃逸到无限远时，我们可以放心地断定，它会把能量从黑洞带走，黑洞的质量必然因此缩减。（这个现象并不违反面积定理，因为量子场的能动张量在视界附近将不满足弱能量条件。）随着质量缩减，表面引力增大，温度也随之升高；这是一个失控的过程，全部质量会在有限的时间内蒸发殆尽。代入数值，可以得到量级如下的寿命

\begin{equation*}
\tau_{\mathrm{BH}} \sim \left( \frac{M}{m_{\mathrm{P}}} \right)^{3} t_{\mathrm{P}} \sim \left( \frac{M}{M_{\odot}} \right)^{3} \times 10^{71}\,\mathrm{sec},
\tag{9.178}
\end{equation*}

其中 $m_{\mathrm{P}} \sim 10^{-5}\,\mathrm{g}$ 是 Planck 质量，$t_{\mathrm{P}} \sim 10^{-43}\,\mathrm{sec}$ 是 Planck 时间。由于 Hubble 时间为 $H_{0}^{-1} \sim 10^{18}\,\mathrm{sec}$，一个太阳质量黑洞的寿命约为宇宙年龄的 $10^{53}$ 倍。这看起来是很长的时间，但我们在这里谈论的是原理性的问题。

你可以看出为什么黑洞熵的问题变得如此严峻：一旦黑洞蒸发殆尽，我们就不能再求助于事件视界来隐藏黑洞那些所谓的态。黑洞不复存在，只剩下它产生的霍金辐射。这种辐射被认为应当是严格热性的（出射粒子中没有隐藏的关联），这一事实意味着，它没有任何办法传递巨量的信息，而指定我们的熵计算所隐含的那些态正需要这些信息。这样一来，如果我们把两个非常不同的初始态分别坍缩成两个质量、电荷和自旋都相同的黑洞，它们就会辐射成两团不可分辨的霍金粒子云。系统在变成黑洞之前、为规定其状态而投入的信息，看来已被抹去；这就是\textbf{信息丢失佯谬}（information loss paradox）。量子场论和广义相对论都以幺正演化为特征——指定早期时刻的一个态所需的信息，与指定较晚时刻的一个态所需的信息严格相等，因为二者通过运动方程联系在一起。但在把 QFT 与 GR 结合的过程中，这种幺正性显然被破坏了。看来很可能我们在论证的某个地方作了某个不恰当的假设，但很难看出在哪里。

传达信息丢失佯谬实质的一个途径，是考察图 9.3 所示的设想中的蒸发黑洞共形图。我们并不真正知道完整的时空应当是什么样子，但这里我们作了似乎合理的假设：形成一个奇点，以及一个与之相伴随的事件视界，二者在黑洞完全蒸发后都消失，留下一个具有 Minkowski 型因果结构的时空。如果我们用 Cauchy 面来思考，问题就显而易见了。一张从类空无限远 $i^{0}$ 延伸到奇点过去侧某个 $r = 0$ 点的非编时面，其未来依赖域将是整个时空，因此这样的面会是一个 Cauchy 面。但是一张延伸到奇点未来侧某个 $r = 0$ 点的类似的面就不是 Cauchy 面了，因为事件视界背后的区域不在它的依赖域之中。因此，由于信息消失于奇点，过去无法由未来反推。换句话说，这个过程看起来是时间不可逆的（在微观意义上，而不仅仅是统计意义上），尽管用来预言它的动力学定律在时间反演下是完全不变的。

%FIG{9.3}: {设想中的蒸发黑洞共形图。能量被霍金辐射带走，黑洞最终完全蒸发殆尽，留下具有 Minkowski 时空因果结构的未来。落过事件视界进入奇点的信息看起来丢失了。}

在讨论信息丢失佯谬时，请记住我们对黑洞蒸发的分析只是在量子场论与广义相对论相耦合的混合理论中进行的，并不是在现实的量子引力理论中。真实世界里究竟可能发生着什么呢？一种可能性是信息真的会丢失，幺正性真的会被破坏，而我们只得学会与之共处。许多物理学家觉得，引入这样一种可预见性的根本失效令人难以接受，而且也有论证表明，幺正性的破坏必然会导致能量守恒的破坏。另一种可能性是，幺正性在我们的世界中看似被破坏，只是因为进入黑洞的信息以某种方式逃逸到了空间一个不相连通的区域——一个婴儿宇宙（baby universe）。广义相对论预言黑洞中心存在奇点，而不是产生一个不相连通的区域；但显然我们正处于量子效应将剧烈改变经典预期的领域，所以我们应当保持开放的心态。

反对信息丢失的一些证据来自弦论。弦论自然地在 10 或 11 维时空中定义，其特征不仅在于一维的延展客体（弦），还在于统称为“膜”（branes）的各种更高维的延展客体。弦论的一个关键方面，是将玻色子与费米子联系起来的高度超对称性。在现实世界中，超对称性如果真的存在，就必须是自发破缺的，因为我们没有观测到与电子具有相同质量和电荷的玻色版本。但作为思想实验的工具，超对称性是无价的。人们可以构造出弦和膜的超对称位形，用来描述各种维度中的黑洞几何。弦论中有一个自由参数（实际上是一个标量场），即弦耦合（string coupling），它控制着引力的强度以及其他各种力的强度。如果我们考虑在某个弦耦合取值下描述黑洞的位形，那么当我们减小耦合时，Schwarzschild 半径最终会缩小到这个位形的尺寸以下，于是该位形就变成一团弱耦合的弦和膜。由于超对称性程度很高，我们可以确信，当我们改变弦耦合时，态的各种特征保持不变；特别是，我们预期自由度的数目（因而熵）不会改变。但在弱耦合区域没有黑洞，我们只有一团由通常自由度组成的“气体”（诚然，是高维中延展客体的气体），它的熵我们应该能够可靠地计算。

Strominger 和 Vafa 对一类特殊的五维超对称黑洞（带有不同种类的荷）考虑了这个过程。\footnote{A. Strominger and C. Vafa, “Microscopic origin of the Bekenstein-Hawking entropy,” \emph{Phys. Lett.} B \textbf{379}, 99 (1996), \texttt{http://arxiv.org/hep-th/9601029}。评述可参见 Johnson (2003)，或 A.W. Peet, “TASI lectures on black holes in string theory,” (2001), \texttt{http://arxiv.org/hep-th/0008241}。}他们得到了一个非凡的结果：系统在弱耦合下的自由度数目，与根据强耦合下的黑洞熵所预言的数目精确相符。由于黑洞熵非平庸地依赖于位形的荷，这种符合看来不大可能是偶然的。后续的研究把这一分析扩展到了其他类型的黑洞，并且不断地发现相符的结果。此外，通过考虑在弱耦合系统上的散射，我们甚至可以计算出预期的黑洞灰体因子；同样，结果与强耦合下的预期相符。因此，至少在弦论中，我们有极好的理由相信，黑洞辐射所隐含的自由度是真实存在的。

遗憾的是，弦论对态的计数，对于“黑洞状态的信息如何以某种方式传给出射的霍金辐射”这个问题，几乎没有提供直接的理解。尽管如此，我们当然应当认真对待事情果真如此的可能性，即使想象这样一个过程究竟如何实际运作还面临着严重的困难。当我们考虑把某些信息——或许以一套百科全书的形式——扔进一个大黑洞，而那时离黑洞蒸发殆尽还早得很的时候，困难就出现了。在这个阶段，黑洞温度很低，表面引力极小，事件视界附近的时空曲率也相当小。从百科全书的角度看，视界处没有任何特别的事情发生，我们应当预期它基本上毫发无损地落过视界。特别是，很难想象百科全书中的信息如何能够传递给早期发射出的霍金辐射。在幺正演化中，信息不可能被复制；它要么随百科全书一起落过视界，要么就必须在视界即将被穿越之前被有效地提取出来，而这似乎不合情理。我们或许可以指望，信息伴随着百科全书进入奇点附近的某个区域，并以某种方式保存在那里，直到黑洞变得很小的晚期。但到那时，大部分辐射粒子已经被发射出去了，最后一阵辐射所能到达的态的数目，一般会小于描述可能已经落入黑洞的各种不同态所需要的数目。

因此，要想认为信息以某种方式编码在出射辐射之中，看来就有必要在很早的时候就把关联编码到霍金粒子之中。我们刚刚论证过这一点很难做到，因为黑洞很大时，视界是一个平淡无奇的地方。摆脱这一困境的一个可以设想的办法，是迈出戏剧性的一步：放弃局域量子场论。换句话说，我们一直隐含地假设，信息可以被合理地描述为定域在空间的某个区域之中；这是通常量子场论的一个无可争议的特征。但量子引力也许有所不同，黑洞所包含的信息或许以某种方式非局域地弥散在视界之上。这个建议本身并不能直接给出把信息送进出射霍金辐射的机制，但它确实使我们给出的、关于为什么这样做会很困难的一些论证受到质疑。

非局域性的一种具体实现，以\textbf{全息原理}（holographic principle）为名。这个思想最初由 't Hooft 和 Susskind 提出：空间一个区域中的自由度数目，不是与该区域的体积成正比（正如在局域场论中所预期的那样），而是与该区域边界的面积成正比。\footnote{评述可参见 R. Bousso, “The Holographic Principle” (2002), \texttt{http://arxiv.org/hep-th/0203101}。}其灵感当然来自黑洞熵，后者按事件视界的面积标度；如果熵计数的是可及态的数目，全息就能解释为什么起作用的是面积，而不是所包围的体积。你也许会担心如何处理闭合宇宙的情形：在其中某个区域可能几乎由整个空间构成，边界却非常小；不过，只要把空间区域换成一组从边界向内延伸的“光片”（lightsheets），就可以表述出一个更具协变性的全息原理版本。全息原理最伟大的胜利是在第 8 章提到过的 AdS/CFT 对应中。在那里，反 de Sitter 背景上的量子引力物理，等价于定义在 AdS 边界上、维数低一维的无引力共形场论。可以设想，我们在宇宙中观测到的所有物理现象，都可以由某个定义在更低维数中的普通非引力理论的非局域全息投影来描述；我们究竟应当如何构造这样一种对应、又如何将它与观测联系起来，目前还完全不清楚，但宇宙学以及宇宙大尺度结构的考虑也许会是一个有希望的切入点。

关于黑洞熵、弦论和全息原理的这些评论，显然无意作为对一个非常活跃的研究领域的细致导引。毋宁说，它们意在指出引力物理前沿正在探索的若干可能性。经典广义相对论是迄今发明的最优美的物理理论，但我们完全有权利期待，GR 与物理学其他领域的综合，将揭示出我们如今只能想象的层层之美。